\exercisegroup

\begin{enumerate}
\def\labelenumi{\arabic{enumi})}
\tightlist
\item
  设 B 为拓扑空间，\((X,p)\) 与 \((Y,q)\) 为 B-空间，\(f\colon X\to Y\) 为满足 \(q\circ f=p\) 的映射。设 A 为拓扑空间，\(g\colon A\to B\) 为连续映射。我们考虑由 f 换基而得的 A-态射 \(f_{A}\colon X_{A}\to Y_{A}\) 。
\end{enumerate}

\begin{enumerate}
\def\labelenumi{\alph{enumi})}
\item
  给出一个例子，其中映射 \(f_{A}\) 连续，映射 g 满射，但映射 f 不连续。
\item
  假设 g 是满射且普遍严格的，且 \(f_{A}\) 连续。证明 f 连续。
\item
  我们假设映射 \(g\) 满射，且映射 \(q\) 与 \(f_{\mathrm{A}}\) 是开的。证明 \(f\) 是开的。
\end{enumerate}

\begin{enumerate}
\def\labelenumi{\arabic{enumi})}
\setcounter{enumi}{1}
\tightlist
\item
  我们重新采用 I, 第 15 页命题 7 的记号。
\end{enumerate}

\begin{enumerate}
\def\labelenumi{\alph{enumi})}
\item
  给出一个例子，其中方块 (12) 与 (14) 是笛卡尔的，但方块 (13) 不是。
\item
  我们假设映射 \(f \circ g\) 是满射且普遍严格的。证明：若方块 (12)、(13) 与 (14) 中有两个是笛卡尔的，则第三个也是。（利用习题 1。）
\end{enumerate}

